\section{Examples}
\subsection{Wishart Libor Model}
The dynamics above show that the Wishart specification provides a very rich structure of the model. In order to get an understanding of the impact of the various parameters we will look first at the impact of $\beta$ followed by the elements $M_{11},M_{22},Q_{11},Q_{22}$ in a model in which all matrices are assumed to be diagonal.
We use the following set of parameters as a benchmark:
\begin{align}
\Sigma_0&=\left(\begin{array}{cc} 3.75&0\\ 0&3.45\end{array}\right)\quad M=\left(\begin{array}{cc} -0.3125*1.0e-003&0\\ 0& -0.5000*1.0e-003\end{array}\right)\nonumber\\
Q&=\left(\begin{array}{cc} 0.034&0\\ 0&0.0420\end{array}\right)\quad \beta=2.7\nonumber
\end{align}
The impact of $\beta$ is quite easy to understand: the process acts by influencing the overall level of the surface. This is due fact that the higher $\beta$ the lower the probability that the process $\Sigma$ approaches 0. It is interesting to note that there is not only a level impact, but also a curvature effect.
\begin{center}
Figure \eqref{fig:1} here
\end{center}
We then may look at the parameters along the diagonals of the matrices $M$ and $Q$. The following claims may be easily checked by looking at the SDE's satisfied by the elements of $\Sigma$. Note that we assumed all eigenvalues of $M$ to lie in the negative real line.. 
\begin{itemize}
\item For $\left|M_{11}\right|\nearrow(\searrow)$ the surface is shifted downwards (upwards).
\item For $\left|M_{22}\right|\nearrow(\searrow)$ the surface is shifted downwards (upwards).
\end{itemize}
The impact is more evident for OTM caplets with short maturities. This is due to the fact that as the process decreases (in matrix sense) the probability that caplets with short maturities are exercised is lowered more than the analogous probability for longer term caplets.
\begin{center}
Figure \eqref{fig:2}  here
\end{center}
We then consider the impact of $Q_{11},Q_{22}$. We have the following:
\begin{itemize}
\item As $Q_{11}\nearrow(\searrow)$ the surface is shifted upwards (downwards). In particular if we multiply $Q_{22}$ by a constant $c>1$, then the increment in the short term is higher for OTM than for ITM caplets. If $c<1$ then the decrease is higher for short term OTM caplets, which is intuitive, given the discussion above.
\item The same impacts, with different magnitudes, is observed also for $Q_{22}$
\end{itemize}
\begin{center}
Figure \eqref{fig:3}  here
\end{center}


\subsubsection{The skew}
Once we fix this model, we are able to compute the so called skew, i.e. the covariation between the Libor rate and its volatility. In order to compute this quantity, we need to derive first the dynamics of the Libor rate in the present model.
The dynamics of the libor rate may be derived along the following steps: we first define the quantity
\begin{align}
B_{T_N-t}(u_k,u_{k+1})=\psi_{T_N-t}(u_k)-\psi_{T_N-t}(u_{k+1})=B_k
\end{align}
\begin{align}
\frac{dL(t,T_k,T_{k+1})}{L(t,T_k,T_{k+1})}=\frac{1+\delta L(t,T_k,T_{k+1})}{L(t,T_k,T_{k+1})}\left((...)dt+Tr\left[B_k(u_k,u_{k+1})d\Sigma_t\right]\right)
\end{align}
In order to preserve analytical tractability, we freeze the coefficients and approximate as follows:
\begin{align}
\frac{1+\delta L(t,T_k,T_{k+1})}{L(t,T_k,T_{k+1})}\approx \frac{1+\delta L(0,T_k,T_{k+1})}{L(0,T_k,T_{k+1})}=C
\end{align}
Then we can write:
\begin{align}
\frac{dL(t,T_k,T_{k+1})}{L(t,T_k,T_{k+1})}&=C\left((...)dt+2\sqrt{Tr\left[QB_k\Sigma B_k^{\top}Q^{\top}\right]}\left(\frac{Tr\left[QB_k\sqrt{\Sigma}dW\right]}{\sqrt{Tr\left[QB_k\Sigma B_k^{\top}Q^{\top}\right]}}\right)\right)\nonumber\\
&=C\left((...)dt+2\sqrt{Tr\left[QB_k\Sigma B_k^{\top}Q^{\top}\right]}d\tilde{W}\right)
\end{align}
The scalar noise driving the factor process may be derived as follows:
\begin{align}
&dTr\left[QB_k\Sigma B_kQ^{\top}\right]=\left(Tr\left[QB_k\beta Q^{\top}QB_kQ^{\top}\right]+2Tr\left[QB_kM\Sigma B_kQ^{\top}\right]\right)dt\nonumber\\
&+2Tr\left[QB_k\sqrt{\Sigma}dWQB_kQ^{\top}\right]\nonumber\\
&=(...)dt+2\sqrt{Tr\left[\Sigma B_kQ^{\top}QB_kQ^{\top}QB_kQ^{\top}QB_k\right]}\frac{Tr\left[QB_kQ^{\top}QB_k\sqrt{\Sigma}dW\right]}{\sqrt{Tr\left[\Sigma B_kQ^{\top}QB_kQ^{\top}QB_kQ^{\top}QB_k \right]}}\nonumber\\
&=(...)dt+2\sqrt{Tr\left[\Sigma B_kQ^{\top}QB_kQ^{\top}QB_kQ^{\top}QB_k\right]}dZ
\end{align}
The covariation between the two noises is calculated as follows:
\begin{align}
\left\langle d\tilde{W},dZ\right\rangle&=\left\langle \frac{Tr\left[QB_k\sqrt{\Sigma}dW\right]}{\sqrt{Tr\left[QB_k\Sigma B_k^{\top}Q^{\top}\right]}},\frac{Tr\left[QB_kQ^{\top}QB_k\sqrt{\Sigma}dW\right]}{\sqrt{Tr\left[\Sigma B_kQ^{\top}QB_kQ^{\top}QB_kQ^{\top}QB_k \right]}} \right\rangle\nonumber\\
&=\frac{\left(\sum{Q_{pq}B_{qr}\sqrt{\Sigma}_{rs}dW_{sp}}\right)\left(\sum{Q_{ab}B_{bc}Q^{\top}_{cd}Q_{de}B_{ef}\sqrt{\Sigma}_{fg}dW_{ga}}\right)}{\sqrt{Tr\left[QB_k\Sigma B_k^{\top}Q^{\top}\right]}\sqrt{Tr\left[\Sigma B_kQ^{\top}QB_kQ^{\top}QB_kQ^{\top}QB_k \right]}}\nonumber\\
&=\frac{\sum{B_{fe}Q^{\top}_{ed}B_{cb}Q^{\top}_{ba}Q_{aq}B_{qr}\sqrt{\Sigma}_{rg}\sqrt{\Sigma}_{gf}}dt}{\sqrt{Tr\left[QB_k\Sigma B_k^{\top}Q^{\top}\right]}\sqrt{Tr\left[\Sigma B_kQ^{\top}QB_kQ^{\top}QB_kQ^{\top}QB_k \right]}}\nonumber\\
&=\frac{Tr\left[B_kQ^{\top}QB_kQ^{\top}QB_k\Sigma\right]dt}{\sqrt{Tr\left[QB_k\Sigma B_k^{\top}Q^{\top}\right]}\sqrt{Tr\left[\Sigma B_kQ^{\top}QB_kQ^{\top}QB_kQ^{\top}QB_k \right]}}\nonumber\\
\end{align}
From this formula we realize that the matrix $Q$ is responsible for the shape of the skew. We also have an indirect impact of $M$ coming from the term $B_k$ which is the difference of two solutions of the Riccati ODE's. 

The presence of $\Sigma$ suggests that in the present framework the skew is stochastic, moreover it can only have positive sign: since we have $\xi_k>\xi_{k+1}$ then we have $u_k\succ u_{k+1}$ and then, according to Lemma 3.2 in \cite{article_Cuchiero} we have that $B_k\in S_d^+$. In all terms in the numerator and the denominator we recognise congruent transformations of matrices in $S_d^+$ which leave the eigenvalues unchanged. The self-duality of $S_d^+$ allows us to claim that all traces are positive, hence we are done.

%\subsubsection{Numerical Tests on the Skew}
%From the dynamics of the single factors we recall that the volatility of $\Sigma_{11}$ is influenced by $Q_{21}$ whereas $\Sigma_{22}$ is influenced by $Q_{12}$.
%We choose the following starting values for the parameters:
%\begin{align}
%\Sigma_0&=\left(\begin{array}{cc} 0.1563&0\\ 0&0.1437\end{array}\right)\quad M=\left(\begin{array}{cc} -0.6944e-03&0\\ 0&-0.1111e-03\end{array}\right)\nonumber\\
%Q&=\left(\begin{array}{cc} 0.0121&0\\ 0&0.0150\end{array}\right)\quad \beta=2.7\nonumber
%\end{align}
%We fix $t=0$ and $\forall k=1,..,N-1$ we consider the skew of the associated libor rate, given by the formula above. We plot the skews for all libor rates, thus obtaining a curve which represent all the skews of the libor rates at time 0. We look at the shape of this curve for different values of $Q_{12},Q_{21}$. In the following plots on the first row we let $Q_{12},Q_{21}$ range from 0 to 0.0036 and let $\Sigma_0$ be diagonal as above. In all cases we observe that the skew for the nearest Libor rates is higher. Increasing $Q_{12}$ or $Q_{21}$ has the same effect: we obtain higher skews. We notice also that the impact of $Q_{21}$ is more pronounced. The explaination is given by the fact that $Q_{21}$ impacts on $\Sigma_{11}$, which is the factor with highest magnitude.
%This happens also if we consider a model in which $\Sigma_0$ is fully populated, for which we report the skews on the second row of the table below, where we set $\Sigma_{12}=\Sigma_{21}=0.6$.
%\begin{center}
%Figure \eqref{fig:7}  here
%\end{center}
%
%We may now look also at the impact of the matrix $M$. When $\Sigma_0$ is diagonal as above, the impact of $M_{12},M_{21}$ is negligible. On the contrary, if we consider where $\Sigma_{12}=\Sigma_{21}=0.6$ as above we have that the skews on the short term are influenced more than those on the long term.
%\begin{center}
%Figure \eqref{fig:8}  here
%\end{center}
%
\subsubsection{The Term Structure of ATM Implied Volatilities}
We proceed to consider the term structure of caplets implied volatilities. When the matrix $\Sigma_0$ is diagonal, the impact of the elements of $Q$ is the same: any increase in the absolute value of any element of $Q$ will result in a steeper term structure of ATM caplet volatilities.
\begin{center}
Figure \eqref{fig:4}  here
\end{center}
Considering a model where $\Sigma_0$ is a full matrix does not influence this result in a significant way.

\subsubsection{More complex adjustments}
In order to appreciate the flexibility of this modelling framework, we focus now on the impact of the off-diagonal elements. We introduce off-diagonal elements in $M$ and $Q$ and look at the relative change in the short term smile (4 months) and the long term smile (32 months).
We introduce a fully populated matrix $\Sigma_0$ and look at the impact of $M_{12}$ and $M_{21}$. Our experiments show that there is a symmetry between the sign of $\Sigma_{0,12}$ and $M_{12}$, $M_{21}$. More precisely, the implied volatility changes are as in table \eqref{tab:ImpliedVolatilityChanges}.	
\begin{table*}[h]
\centering
	\begin{tabular}{ccc}
	 &$\Sigma_{0,12}>0$&$\Sigma_{0,12}<0$\\
	 $M_{12}>0$&Increase&Decrease\\
	 $M_{12}<0$&Decrease&Increase \\
	 &&\\
   $M_{21}>0$&Increase&Decrease\\
	 $M_{21}<0$&Decrease&Increase \\
	\end{tabular}
	\caption{Implied volatility changes: relation between $\Sigma_{0,12}$ and $M_{12},M_{21}$}
	\label{tab:ImpliedVolatilityChanges}
\end{table*}
The reason for this symmetry is to be looked for in the drift part of the dynamics of the single elements of the matrix process $\Sigma$.
Next we look at the impact of $Q_{12}$ and $Q_{21}$. To this end we model $Q$ as a symmetric matrix and set $Q_{21}=Q_{12}=\rho\sqrt{Q_{11}Q_{22}}$ for a real parameter $\rho$. Also in this case we we recognize two main shapes of the adjustment that we denote by $A,B$.

\begin{table*}[h]
\centering
	\begin{tabular}{ccc}
	 &$\Sigma_{0,12}>0$&$\Sigma_{0,12}<0$\\
	 $\rho>0$&A&B\\
	 $\rho<0$&B&A\\
	\end{tabular}
	\caption{Implied volatility changes: relation between $\Sigma_{0,12}$ and $\rho$}
	\label{tab:ImpliedVolatilityChanges_rho}
\end{table*}

We now proceed to perform other numerical tests which will show that our modelling framework has a certain degree of flexibility. For these tests we set:
\begin{align}
M&=\left(\begin{array}{cc} -0.3125*1.0e-003&0\\ 0& -0.5000*1.0e-003\end{array}\right)\nonumber\\
Q&=\left(\begin{array}{cc} 0.02&\rho \sqrt{Q_{11}Q_{22}}\\ \rho \sqrt{Q_{11}Q_{22}}&0.0420\end{array}\right)\quad \beta=2.7\nonumber
\end{align}
So basically $M$ is parametrized  as before but $Q$ is symmetric and equiped with a parameter $\rho$ which summarizes the information on the off-diagonal elements. We require $\Sigma_0=\Sigma_\infty$, where $\Sigma_\infty$ is given by the solution of the Lyapunov equation \eqref{Lyap}. After that we perturbate $\Sigma_0$ in order to include off-diagonal elements and set $\Sigma_{0,12}=\Sigma_{0,21}=2$.
We have a good degree of control on the term structure of ATM implied volatilities. In particular, we may have larger percentage shifts in the long-term w.r.t. the short-term ATM implied volatility, or, for $\rho=-0.6$ we may even reproduce a situation where the short term ATM implied volatility increases whereas the long-term ATM implied volatility decreases.
\begin{center}
Figure \eqref{fig:5}  here
\end{center}

If we adopt the same kind of parametrization for the matrix $M$ by introducing a second parameter $\rho_2$ than we have further flexibility because we can impose many different combinations of $\rho,\rho_2$. For example, the image below shows that we are in practice able  to isolate an effect on the term structure of ATM implied volatility: we have a moderate change for ITM caplets, OTM caplets are practically unchanged, but the shape of the term structure of ATM implied volatility is modified in a significant way.
\begin{center}
Figure \eqref{fig:9}  here
\end{center}

Finally, just for illustrative porpuses we report a prototypical Caplet volatility surface generated by the model.
\begin{center}
Figure \eqref{fig:7}  here
\end{center}
As far as Swaptions are concerned an example of ATM implied volatility surface for different expiries and underlying swap lengths is given below.
\begin{center}
Figure \eqref{fig:8}  here
\end{center}
